\section{Restarting FHI-aims calculations}
\label{Sec:restarts}
\emph{This section gives background information on restarts in FHI-aims.}

\subsection{General restart procedure}
Calculations can be restarted from previous wavefunctions using the keyword
\keyword{restart}. In general, restarting the same calculation on the same
computer using the same number of MPI-tasks will always work.

\paragraph{Important --- Rotated systems:} FHI-aims wavefunctions are not
symmetry-aware. While the same restart file can be used for molecules
\emph{translated} by a vector, this does not work for \emph{rotated} molecules.
For a way to restart systems with rotated geometries or superpositions of
molecular densities see Section~\ref{Sec:rotated-restarts}.

FHI-aims has two main means of saving the wavefunction for a later restart,
depending on the used keywords.

\paragraph{Kohn--Sham eigenvector (\emph{"wavefunction"}):} This is the default
case always used for calculations using the LAPACK solver. The saved file
contains some information on the system (number of basis functions, states,
spin, k-points), the Kohn--Sham eigenvector, eigenvalues and occupations.

\paragraph{Kohn--Sham density matrix (\emph{"density matrix"}):} This variant is
only used for calculations using the scaLAPACK or ELPA solver. The saved file
contains the density matrix ($\hat{n}_{ij}=\sum_l f_l c_{il} c_{jl}$), either in
full or sparse storage.

The corresponding keywords determining which version of storage is used are
\keyword{KS\_method} and \keyword{density\_update\_method}. See
Fig.~\ref{fig:restart_scheme} for possible options and outcomes.

\begin{figure}[hb]
  \centering
  \includegraphics[width=0.9\textwidth]{./imgs/restart_scheme}
  \caption{Scheme of different restart types depending on used methods.}
  \label{fig:restart_scheme}
\end{figure}

For technical reasons there is not one single way, but the above mentioned
variants with again, variants.

The \emph{density matrix} based restart always
writes one file per spin-channel, no matter how many parallel MPI tasks are
used.
The \emph{wavefunction} based restart generally writes a single restart file
for all cluster calculations and one file per MPI-task for periodic
calculations. Due to this periodic calculations using the \emph{wavefunction}
restart need to be restarted with the same number of MPI-tasks.

\subsection{Mixing variants - the "force\_single\_restartfile" option}
For a subset of possible calculations the above described differences are not
valid and it is possible to restart different calculations using one type of
restart file. This is the case for all possible cluster calculations (using
either the LAPACK or scaLAPACK/elpa solver) and for periodic calculations with
only 1 k-point, $\Gamma$-only (using either the LAPACK or scaLAPACK/elpa solver).

This is implemented via the keyword \keyword{force\_single\_restartfile}. Using
this functionality, it is for example possible to profit from the parallel ELPA
solver for huge systems and still get a file containing the KS-wavefunction for
post-processing.

\subsection{Comments on the 'restart' starting point and on self-consistency}
As noted in the description of the \keyword{restart} keyword, 
it is important to note that the \keyword{restart} infrastructure
corresponds to a restart from the last Kohn-Sham orbitals, not from
the last density. 

In practice, this means that the code will restart
from the last \emph{unmixed} Kohn-Sham density, not from the last
\emph{mixed} density. When restarting from a non-selfconsistent
starting point, this can lead to unexpected jumps in the calculated
non-selfconsistent total energy between the ``old'' and the ``new'' 
(restarted) FHI-aims run.

Only the self-consistent total energy is truly meaningful. This quantity (the
self-consistent) total energy should be the same for the same
stationary density, when approached from different starting densities.
However, note additionally that some systems may exhibit several
different self-consistent stationary densities even for 
the exact same atomic positions and for the exact same density functional. 
A simple example are
antiferromagnetic vs. ferromagnetic spin states in some systems. In
such cases, the true ground state in a DFT sense is the stationary density 
that yields the lowest energy. It can be found by way of a global search for
different stationary densities, usually by varying the initial density guess.


\subsection{Rotating the FHI-aims wavefunction}
\emph{This feature is new and tested for the case of weakly interacting organic
        molecules as can be found in organic crystals. If you encounter any
        problems or strange behaviour, please let us know!}

In many molecular systems a rotation does not change the electron density
(or wavefunction) of a system. Therefore, the same wavefunction can be used
again to restart such a calculaction.

This is especially true for weakly interacting systems such as van-der-Waals
bound organic crystals. In these cases, the wavefunction of each single molecule
in the crystals is not too different from the isolated wavefunction of a single
molecule. This can be exploited to start a FHI-aims calculation of huge
molecular systems from a \emph{superposition of molecular densities} instead of
the usual \emph{superposition of atomic densities}. 

This functionality is implemented as an external tool using Python and can be
found in the FHI-aims repository itself (Python-package \emph{aimsutils}) or at
\url{https://gitlab.lrz.de/theochem/aimsutils}.

If you use this functionality, please cite \emph{C. Schober, K. Reuter, H.
        Oberhofer, J. Chem. Phys., 144:054103, 2016}\cite{Schober2015}.

\subsubsection*{Important pre-requisites}
This functionality only works for restart files containing the wavefunction of
the system, not the density matrix. Therefore, please make sure you have
wavefunction restart files or use the option
\keyword{force\_single\_restartfile}. The FHI-aims output file of the
calculation is also necessary.

Rotations can be defined via Euler angles in \emph{zyz} convention or quaternions. 
To avoid any confusions with rotational conventions (especially with Euler
angles), the input geometry will be rotated and saved using the same rotation
matrix. Please be sure to check if your expected and the actual rotation are the
same (or directly use the generated \texttt{geometry.in}). 

The final rotated restart file is again in the wavefunction format.

\begin{verbatim}
- example_calc
|
+-- control.in              
+-- geometry.in             # necessary 
+-- aims.out                # necessary 
+-- restart.wavefunction    # name can be anything, will be parsed from outfile
\end{verbatim}

\subsubsection*{Examples}
The two examples are shown using pre-defined scripts, but it is
totally possibly to use the Python package directly. The full Python API is
available via the Sphinx-documentation of the \emph{aimsutils} package.

\paragraph{Rotation of a single system:} The script \texttt{RotateSingle.py} can be
used to create the rotated geometry and restart for a single calculation.
\begin{verbatim}
        $: # RotateSingle alpha beta gamma -t x y z
                          '-----.-----'    '---.---'
                          Euler angles   translation vector
        $: RotateSingle 45 0 123 5.4 2.3 6.3

        $: # RotateSingle x_i x_j x_k x_l -t x y z
                          '------.------' '---.---'
                            quaternion   translation vector
        $: RotateSingle 0.33 -0.12 0.44 0.84 5.4 2.3 6.3
\end{verbatim}
This will create a subfolder \emph{rotated} with the new \texttt{restart.combined}
and \texttt{geometry.in}.

\paragraph{Rotation and combination of wavefunctions} The script
\texttt{RotateMany.py} reads a instruction file (\texttt{rotations.in}) with the
following format:
\begin{verbatim}
            # Comments or empty lines are ignored

            # title will also be foldername
            title project_x
            # lattice vector for new cell if any
            lattice_vector 11.2 0.00 0.00
            lattice_vector 0.00 10.5 0.00
            lattice_vector 0.00 0.00 10.7
            # now all rotated molecules
            # Euler angles (zyz)
            # alpha beta gamma x y z parent_molecule
            33.4 12.4 167.0   10.  0.  0.  calc_2
             0.  45.    0.     3.5 2.1 5.4 calc_1
           130.  90.   23.     0.  2.6 12. calc_3
            # OR Quaternion
            # x_i x_j x_k  x_l   x y z
            -0.4 0.15 0.33 -0.84 4 0 0 calc_2
            -0.2 0.45 0.13  0.84 0 2 0 calc_1
             0.0 0.00 0.00 -0.32 4 4 5 calc_3
\end{verbatim}

The necessary converged FHI-aims calculations must be in subfolders with the
appropriate name:
\begin{verbatim}
- calc_something/ (main calculation folder)
|
+-- rotations.in
+-- calc1/
    |
    +-- ...
+-- calc2/
    |
    +-- ...
+-- calc3/
    |
    +-- ...

\end{verbatim}
This will create a subfolder \texttt{project\_x} with the rotated wavefunction
assembled from the individual calculations defined in the script
(\texttt{restart.combined}) and geometry (\texttt{geometry.in}).

\subsubsection*{Theory}
\label{Sec:rotated-restarts}
The following should give a summary of how the rotation of wavefunctions can be
done with FHI-aims. For details on the used methods and mathematics, please
have a look at the references.

In FHIaims, a basis function is defined as
\begin{align}
    \Phi_{i,lm} = \frac{u_i(r)}{r}\cdot S_{l,m}(\theta, \phi),
\end{align}
with $S_{l,m}(\theta, \phi)$ being real valued spherical harmonics. These are
obtained from the complex spherical harmonics $Y_{lm}$ via
\begin{align}
        S_{l,m}(\theta, \phi) = 
    \begin{cases} 
            \frac{(-1)^m}{\sqrt(2)}(Y_{lm} + Y_{lm}^*) & m > 0\\
            Y_{l0} & m=0\\
            \frac{(-1)^m}{i\sqrt(2)}(Y_{l|m|} - Y_{l|m|}^*) & m < 0\\
    \end{cases}
\end{align}
A wavefunction is then given by
\begin{align}
        \Psi_k(r) = \sum_{i=1}^{\text{n\_basis}}c_i^k\Phi_i(r)
\end{align}
with the coefficients $c_i^k$. 

A rotation of the molecule (rotation matrix $\mathbf{R}$ leads to the same set
of YLMs $\mathcal{Y}$ (as they are fixed with respect to the xyz-coordinate
system), but with different coefficients $\mathbf{c}$ (new linear combination
of basis functions):

\begin{align}
    \mathbf{c'} = \mathbf{R}\mathbf{c}
\end{align}

The rotation of complex YLMs can be done with Wigner~D matrices ($D^l_{mm'}$).
For real YLMs different schemes are
available\cite{lessig_efficient_2012,aubert_alternative_2013,blanco_evaluation_1997}.
Due to the non-default YLM convention in FHIaims we construct our rotation
matrices for different $l$ from the complex Wigner~D matrices via the
transformation matrix $\mathbf{C}^l$ (reference:
\cite{blanco_evaluation_1997}):
\begin{align}
        S_{l,m} = \mathbf{C^l}Y_{l,m}
\end{align}
The matrix $\mathbf{C}$ is constructed according to
\cite{blanco_evaluation_1997} with the constraint of the different sign
convention in FHIaims. The real rotation matrix $\Delta^l(R)$,
\begin{align}
        \Delta^l(R) = \left(\mathbf{C}^l\right)^*\mathbf{D}^l(R)\left(\mathbf{C}^l\right)^t,
\end{align}
is then used to obtain the rotated coefficients for each $l, m$ for each basis
function in the system.
\begin{align}
    \mathbf{c}'^{l} = \Delta^l(R)\mathbf{c}^{l}
\end{align}
